\section{What does an algorithm cost to propose?}
\label{sec:accounting}

Token yield is meaningful only after specifying the object, unit and phase.
Four quantities should remain distinct: retained model output, compiler-qualified
source, expansion from an internal representation, and useful algorithms per
unit resource. A deterministic emitter can retain every printed token while
producing an unhelpful procedure. A 512-coordinate vector is not 512 language-model
tokens, and source expanded from it is not evidence of near-100\% discovery
efficiency.

For directly observable text under one pinned analysis tokenizer, let $S$ be
source spans retained from model output and $O$ all visible output, including
discarded drafts and tool-call text. Let $R$ be additional, nonoverlapping
reasoning tokens and $I$ all processed input in compatible units. Then
\begin{equation}
 Y_{\mathrm{visible}}=\frac{S}{O},\qquad
 Y_{\mathrm{all}}=\frac{S}{I+O+R}.
\end{equation}
The second ratio is unavailable when hidden reasoning cannot be counted in the
same units. Provider-native usage is a parallel ledger, preserving raw overlap
semantics: reasoning may already be included in output total, and cached input
may already be a subset of input total. Neither is added twice. Deduplicated
input content is a separate analysis, not processed or billable usage. Missing
counters remain unknown rather than zero.

For a latent-to-source pipeline, report learned construction decisions and
emitted source bytes or reference tokens separately. Their ratio is an
\emph{expansion factor}, potentially greater than one. Repository lexical source
counts and compile-yield counters retain their own definitions; they are not
frontier model-token measurements \citep{jope}. Count every attempted output,
including refusals and failures, before reporting useful-source yield.

Historical effort records illuminate these distinctions but do not provide a
controlled model ranking. Astra and Sol estimates use attachment-availability
clock proxies and unmetered token ranges. Report C is repository-attributed to
Opus: its later 94k-plus output already includes 27k reasoning, and belongs to
post-freeze admission/integration. The anonymous single-shot report D remains
unresolved; its source dimensions do not match the located Gemini replacement.
The historical-effort supplement preserves these scopes. This paper performs no new
provider metering or source-authoring experiment.

Measure monotonic phase clocks for parse/preprocess, encode, decode, lower,
commit, compile, fit and predict. Report cold and warm latency, concurrency,
hardware, precision, completion and tail latency. Source tokens divided by a
whole authoring session measure artifact delivery, not neural streaming speed.
Timeouts are censored outcomes, not missing slow cases. A cross-provider service
comparison may have unknown hardware; it cannot establish matched-FLOP efficiency.

For $N$ deployments, the economic comparison is
\begin{align}
 C_{\mathrm{IZ}}(N)&=C_{\mathrm{history}}+C_{\mathrm{train}}+N c_{\mathrm{IZ}},\\
 C_{\mathrm{U}}(N)&=C_{\mathrm{author}}+N c_{\mathrm{U}},
\end{align}
where each deployment cost includes encoding or adaptation, compilation/cache,
fitting and prediction as applicable. A conditional frontier author pays
generation per context; a universal source author does not. Charge all attempts,
and disclose unknown upstream pretraining costs symmetrically. A positive
break-even requires matched quality and positive per-deployment savings. The
operational endpoint is therefore complete cost or latency to a useful fitted
predictor, not the amount of code printed per second.
